\section{课程定位与核心命题}

本讲由 Ohio State University 的 Yu Su 教授主讲，主题是语言智能体（Language Agents）的三个关键能力：Reasoning、Memory 与 Planning。授课并没有把 Agent 当作单一产品形态，而是将其作为一个可扩展的系统对象：模型需要在动态环境中持续感知、决策、执行、复盘和改进。

\begin{importantbox}{本讲核心主张}
\textit{``Reasoning is for better acting.''} 推理不是为了在纸面题目上展示步骤，而是为了在真实环境中做出更好的动作选择。对应地，Memory 不是日志堆积，Planning 也不是静态步骤列表；它们都服务于长期任务成功率与稳定性。
\end{importantbox}

Yu Su 先回顾了过去两年 Agent 热潮中常见的叙事偏差：一方面，外界对 ``2025 是 Agent 元年'' 的预期极高；另一方面，团队在生产实践中经常面对成本失控、长链路任务崩溃和评测口径不稳定等问题。课程的价值就在于用工程化框架替代口号化叙事。

\begin{warningbox}{高热度不等于高可用}
\begin{itemize}
    \item Demo 成功不代表任务分布下的稳定成功；
    \item 单次推理质量高不代表多轮执行质量高；
    \item 模型分数领先不代表系统总成本可接受。
\end{itemize}
\end{warningbox}

\subsection{本章小结}
本章明确了整讲的判断标准：Agent 研究应以系统行为质量为中心，而不是以孤立模型能力为中心。后续关于 Reasoning、Memory、Planning 的讨论，均围绕 ``如何提升真实任务成功率'' 展开。

